%%%%%%%%%%%%%%%%%%%%%%%%%%%%%%%%%%%%%%%%%%%%
%                 PACKAGES
%%%%%%%%%%%%%%%%%%%%%%%%%%%%%%%%%%%%%%%%%%%%

% LANGUAGE AND ENCODING
\usepackage[utf8]{inputenc}      % Direct UTF-8 character input
\usepackage[english]{babel}      % Language-specific hyphenation
\usepackage{xstring}

% CORE MATHEMATICS
\usepackage{amsmath}             % AMS math environments (align, gather, etc.)
\usepackage{amssymb}             % AMS math symbols
\usepackage{amsthm}              % Theorem environments
\usepackage{amsfonts}            % AMS fonts
\usepackage{mathtools}           % Extensions to amsmath
\usepackage{esint}               % Extended integral symbols
\usepackage{cancel}              % Diagonal cancellation in math
\usepackage{tcolorbox}           % Good boxes
\tcbuselibrary{most}

\usepackage{lmodern}
\usepackage{bm}                  % Bold math symbols (Greek, etc.)
\usepackage{mathrsfs}            % Ralph Smith's script font
\usepackage{tensor}              % Tensor index notation

% UNITS AND QUANTITIES
\usepackage{siunitx}             % Professional number and unit formatting
\NewCommandCopy\uqty\qty         % Save siunitx's \qty as \uqty
\usepackage{silence}
\WarningFilter{siunitx}{Detected the "physics" package}
\ExplSyntaxOn
\msg_redirect_name:nnn { siunitx } { physics-pkg } { none }
\ExplSyntaxOff

% TABLES AND BOXES
\usepackage{tabulary}            % Tables with automatic width adjustment
\usepackage{mdframed}            % Framed boxes with backgrounds
% \usepackage{chemformula}       % Alternative chemistry notation (commented)

% MATH TOOLS
\usepackage{centernot}           % Centered negation symbol
\usepackage{dcolumn}             % Decimal alignment in tables
\usepackage[shortlabels]{enumitem} % Customizable lists

% GRAPHICS AND FIGURES
\usepackage{graphicx}            % Include images
\usepackage{float}               % Force figure placement with [H]
\usepackage{tikz}                % Powerful diagram drawing
\usetikzlibrary{positioning,arrows,trees,3d,angles} % TikZ libraries
\usepackage{pgfplots}            % Scientific plots based on tikz
\pgfplotsset{compat=1.18}        % PGFPlots compatibility mode
\usepackage{subcaption}          % Subfigures (a), (b), etc.
\usepackage{caption}             % Customize caption formatting
\usepackage{booktabs}            % Professional tables with \toprule

% MULTI-COLUMN AND ARRAYS
\usepackage{multicol}            % Multi-column text layout
\usepackage{array}               % Enhanced table column types

% CHEMISTRY
\usepackage{chemfig}             % Draw molecular structures
\usepackage[version=4]{mhchem}   % Chemical equations in text (\ce{H2O})

% CIRCUITS
\usepackage{circuitikz}          % Circuit diagrams with tikz

% TEXT FORMATTING
\usepackage{blindtext}           % Lorem ipsum placeholder text
\usepackage{titlesec}            % Customize section headings
\usepackage{fancyhdr}            % Custom headers and footers
\usepackage{csquotes}            % Context-sensitive quotes

% HYPERLINKS
\usepackage[colorlinks=true, allcolors=blue]{hyperref} % Clickable links

% PAGE GEOMETRY
\usepackage[top=1.0in,bottom=1.0in,left=1.0in,right=1.0in,marginparwidth=1.75cm]{geometry}

% CODE LISTINGS
\usepackage{xcolor}              % Color definitions
\usepackage{listings}            % Code syntax highlighting
\usepackage{xparse}              % Robust command definitions
\usepackage{braket}              % Dirac bra-ket notation

% PHYSICS PACKAGES (LOADED LAST)
\usepackage[italicdiff]{physics} % Original physics package (d in italics)
\usepackage{physics2}            % Modern physics package
\usephysicsmodule{ab}            % Auto-braces module
\usephysicsmodule{ab.legacy}     % Legacy \abs, \norm, \eval
\usephysicsmodule{nabla.legacy}  % \grad, \div, \curl operators
\usephysicsmodule{op.legacy}     % Operators like \tr, \Re, \Im
\usephysicsmodule{diagmat, xmat} % Diagonal and special matrices
\usephysicsmodule{braket, ab.braket} % Dirac notation
\let\lpc\laplacian               % Alias for Laplacian

%%%%%%%%%%%%%%%%%%%%%%%%%%%%%%%%%%%%%%%%%%%%
%             NEW COMMANDS
%%%%%%%%%%%%%%%%%%%%%%%%%%%%%%%%%%%%%%%%%%%%

%%% MATH OPERATORS
\DeclareMathOperator{\sinc}{sinc}
\DeclareMathOperator{\sgn}{sgn}

%%% INVERSE TRIGONOMETRIC OPERATORS
\renewcommand{\asin}{\sin^{-1}}
\renewcommand{\acos}{\cos^{-1}}
\renewcommand{\atan}{\tan^{-1}}
\renewcommand{\acsc}{\csc^{-1}}
\renewcommand{\asec}{\sec^{-1}}
\renewcommand{\acot}{\cot^{-1}}

\newcommand{\asinh}{\sinh^{-1}}
\newcommand{\acosh}{\cosh^{-1}}
\newcommand{\atanh}{\tanh^{-1}}
\newcommand{\acsch}{\csch^{-1}}
\newcommand{\asech}{\sech^{-1}}
\newcommand{\acoth}{\coth^{-1}}
% \DeclareMathOperator{\arcsin}{arcsin}
% \DeclareMathOperator{\arccos}{arccos}
% \DeclareMathOperator{\arctan}{arctan}
% \DeclareMathOperator{\arccsc}{arccsc}
% \DeclareMathOperator{\arcsec}{arcsec}
% \DeclareMathOperator{\arccot}{arccot}
\DeclareMathOperator{\arsinh}{arsinh}
\DeclareMathOperator{\arcosh}{arcosh}
\DeclareMathOperator{\artanh}{artanh}
\DeclareMathOperator{\arcsch}{arcsch}
\DeclareMathOperator{\arsech}{arsech}
\DeclareMathOperator{\arcoth}{arcoth}

%%% NUMBER SETS
\newcommand{\R}{\mathbb{R}}          % Real numbers
\newcommand{\Z}{\mathbb{Z}}          % Integers
\newcommand{\Q}{\mathbb{Q}}          % Rational numbers
\newcommand{\C}{\mathbb{C}}          % Complex numbers
\newcommand{\N}{\mathbb{N}}          % Natural numbers
\newcommand{\K}{\mathbb{K}}          % Field R or C
\newcommand{\bb}[1]{\mathbb{#1}}     % Generic blackboard bold

%%% PARTIAL DIFFERENTIALS
\newcommand{\p}{\partial}            % Partial symbol

%%% VECTOR CALCULUS
\newcommand{\krondelta}[1]{\delta_{#1}}
\newcommand{\levicivita}[1]{\varepsilon_{#1}}
\newcommand{\ihat}{\vu*{\imath}}
\newcommand{\jhat}{\vu*{\jmath}}
\newcommand{\khat}{\vu*{k}}
\newcommand{\rhohat}{\vu*{\rho}}
\newcommand{\phihat}{\vu*{\phi}}
\newcommand{\thetahat}{\vu*{\theta}}
\newcommand{\rhat}{\vu*{r}}
\newcommand{\inv}{^{-1}}
\newcommand{\vell}{\boldsymbol{\ell}}

%%% DIFFERENTIALS AND TRANSFORMS
\newcommand{\laplace}[1]{\mathscr{L}\qty{#1}}
\newcommand{\invlaplace}[1]{\mathscr{L}^{-1}\qty{#1}}
\newcommand{\unitstep}[1]{\mathscr{U}\qty{#1}}
\newcommand{\intfty}{\int_{-\infty}^{+\infty}}

%%% PARENTHESES
\newcommand{\floor}[1]{\lfloor #1 \rfloor}
\newcommand{\ceiling}[1]{\lceil #1 \rceil}
\newcommand{\round}[1]{\lfloor #1 \rceil}
\newcommand{\nvec}[3]{\left\langle #1 , #2 , #3 \right\rangle}
\newcommand{\tline}[1]{\hrule \vspace{#1} \hrule}
\newcommand{\blank}[1]{\hspace*{#1}}

%%% TRANSFORM OPERATORS
\newcommand{\Laplace}{\mathscr{L}}
\newcommand{\Fourier}{\mathscr{F}}

%%% LOGICAL OPERATORS
\newcommand{\union}{\cup}
\newcommand{\intersect}{\cap}
\newcommand{\defined}{:=}

%%% COMPLEX ANALYSIS
\newcommand{\Arg}{\operatorname{Arg}}
\newcommand{\Hol}{\operatorname{Hol}}
\newcommand{\Log}{\operatorname{Log}}

%%% LINEAR ALGEBRA
\newcommand{\Span}[1]{\operatorname{span}\qty(#1)}
\newcommand{\Rank}[1]{\operatorname{rank}\qty(#1)}
\newcommand{\nullity}[1]{\operatorname{nullity}\qty(#1)}
\DeclareMathOperator{\image}{im}

%%% TOPOLOGY
\newcommand{\Id}{\operatorname{Id}}
\newcommand{\Image}[1]{\operatorname{Im}\qty(#1)}
\newcommand{\comp}{^{c}}
\newcommand{\Prob}[1]{\operatorname{Pr}\qty(#1)}
\newcommand{\du}{\mathrm{d}}
\newcommand{\topo}[1]{\mathcal{#1}}

%%% MISCELLANEOUS
\newcommand{\notimplies}{\centernot\implies}

\NewDocumentCommand{\dimu}{O{\mkern-3mu} m}{
\left[
#1
\left[#2\right]
#1
\right]
}

\newcommand{\id}{\operatorname{id}}
\newcommand{\contradiction}{\Rightarrow\!\Leftarrow}
\newcommand{\transition}[3]{{_{#2}\mathbf{#1}_{#3}}}
\newcommand{\barline}{\vspace{10pt}\hrule\vspace{10pt}}
\newcommand{\eqques}{\stackrel{?}{=}}
\newcommand{\Gammaf}[1]{\Gamma\qty(#1)}
\renewcommand\qedsymbol{$\square$}


%%% Matrices
\newcommand{\pmat}[1]{\qty(\mqty{#1})}
\newcommand{\bmat}[1]{\qty[\mqty{#1}]}
\newcommand{\vmat}[1]{\qty{\mqty{#1}}}
\newcommand{\dmat}[1]{\abs{\mqty{#1}}}

%%%%%%%%%%%%%%%%%%%%%%%%%%%%%%%%%%%%%%%%%%%%
%       CUSTOM HEADER / SECTION STYLES
%%%%%%%%%%%%%%%%%%%%%%%%%%%%%%%%%%%%%%%%%%%%

\fancypagestyle{titleheader}{
  \fancyhead{}
  \fancyhead{\vspace{2cm}\tline{3pt}\fancyfoot[c]{\thepage}}
  \renewcommand{\headrulewidth}{0pt}
}
\newcommand{\titlebars}{
  \maketitle
  \tline{3pt}
  \thispagestyle{titleheader}
}

%%% CUSTOM SUBSUBSUBSECTION (4-LEVEL HEADINGS)
\titleclass{\subsubsubsection}{straight}[\subsection]
\newcounter{subsubsubsection}[subsubsection]
\renewcommand\thesubsubsubsection{\thesubsubsection.\arabic{subsubsubsection}}
\renewcommand\theparagraph{\thesubsubsubsection.\arabic{paragraph}}
\titleformat{\subsubsubsection}{\normalfont\normalsize\bfseries}{\thesubsubsubsection}{1em}{}
\titlespacing*{\subsubsubsection}{0pt}{3.25ex plus 1ex minus .2ex}{1.5ex plus .2ex}

\makeatletter
\renewcommand\paragraph{\@startsection{paragraph}{5}{\z@}{3.25ex \@plus1ex \@minus.2ex}{-1em}{\normalfont\normalsize\bfseries}}
\renewcommand\subparagraph{\@startsection{subparagraph}{6}{\parindent}{3.25ex \@plus1ex \@minus .2ex}{-1em}{\normalfont\normalsize\bfseries}}
\def\toclevel@subsubsubsection{4}
\def\toclevel@paragraph{5}
\def\toclevel@subparagraph{6}
\def\l@subsubsubsection{\@dottedtocline{4}{7em}{4em}}
\def\l@paragraph{\@dottedtocline{5}{10em}{5em}}
\def\l@subparagraph{\@dottedtocline{6}{14em}{6em}}
\makeatother
\setcounter{secnumdepth}{4}
\setcounter{tocdepth}{4}

%%% CUSTOM SECTION TITLE FONT SIZES
\titleformat*{\section}{\LARGE\bfseries}
\titleformat*{\subsection}{\Large\bfseries}
\titleformat*{\subsubsection}{\large\bfseries}
\setlength{\parskip}{10pt}
%%%%%%%%%%%%%%%%%%%%%%%%%%%%%%%%%%%%%%%%%%%%
%             ENVIRONMENTS
%%%%%%%%%%%%%%%%%%%%%%%%%%%%%%%%%%%%%%%%%%%%

\newenvironment{amatrix}[1]{%
  \left(
  \begin{array}{@{}*{#1}{c}|c@{}}}{
\end{array}\right)}

\newenvironment{abmatrix}[1]{%
  \left[
  \begin{array}{@{}*{#1}{c}|c@{}}}{
\end{array}\right]}

\newenvironment{problem}{
  \hrule
  \vspace{5pt}
}{
  \vspace{10pt}\hrule\vspace{10pt}
}

\newenvironment{solution_pset}{
  \noindent\textbf{\Large Solution}
}{
  \newpage
}

%%%%%%%%%%%%%%%%%%%%%%%%%%%%%%%%%%%%%%%%%%%%
%          THEOREM STYLES
%%%%%%%%%%%%%%%%%%%%%%%%%%%%%%%%%%%%%%%%%%%%

% \theoremstyle{definition}
% \newtheorem{theorem}{Theorem}[subsection]
% \newtheorem{corollary}{Corollary}[theorem]
% \newtheorem{lemma}[theorem]{Lemma}
% \newtheorem{definition}{Definition}[subsection]
% \newtheorem{example}{Example}[subsection]

% \theoremstyle{remark}
% \newtheorem*{remark}{Remark}
% \newtheorem*{solution}{Solution}
% \newtheorem*{note}{Note}

%%%%%%%%%%%%%%%%%%%%%%%%%%%%%%%%%%%%%%%%%%%%
%        PYTHON LISTING STYLE
%%%%%%%%%%%%%%%%%%%%%%%%%%%%%%%%%%%%%%%%%%%%

\definecolor{vscodeBackground}{RGB}{255,255,255}
\definecolor{vscodeForeground}{RGB}{10,10,10}
\definecolor{vscodeKeyword}{RGB}{86,156,214}
\definecolor{vscodeString}{RGB}{206,145,120}
\definecolor{vscodeComment}{RGB}{106,153,85}
\definecolor{vscodeFunction}{RGB}{220,220,170}
\definecolor{vscodeOperator}{RGB}{212,212,212}
\definecolor{vscodeNumber}{RGB}{181,206,168}

\lstdefinestyle{vscodepython}{
  backgroundcolor=\color{vscodeBackground},
  basicstyle=\ttfamily\small\color{vscodeForeground},
  commentstyle=\color{vscodeComment}\itshape,
  keywordstyle=\color{vscodeKeyword}\bfseries,
  stringstyle=\color{vscodeString},
  numberstyle=\tiny\color{vscodeForeground},
  identifierstyle=\color{vscodeForeground},
  morekeywords={True, False, None},
  frame=single,
  rulecolor=\color{vscodeForeground},
  breaklines=true,
  numbers=left,
  numbersep=5pt,
  showstringspaces=false,
  tabsize=4,
  captionpos=b,
}
\lstset{style=vscodepython, language=Python}


%% EMPHASIS BOX

\newtcolorbox{emphasis}{
  colback=white,        % White background
  colframe=black,       % Black border
  boxrule=0.5pt,        % Border thickness
  arc=0pt,              % Sharp corners (0pt = no rounding)
  left=6pt,             % Left padding
  right=6pt,            % Right padding
  top=6pt,              % Top padding
  bottom=6pt,           % Bottom padding
  boxsep=0pt,           % No extra separation
  before upper={
    \setlength{\abovedisplayskip}{0pt}
    \setlength{\abovedisplayshortskip}{0pt}
    \setlength{\belowdisplayskip}{0pt}
    \setlength{\belowdisplayshortskip}{0pt}
  }
}
